\chapter*{Preface}
\addcontentsline{toc}{chapter}{Preface}

Every system that survives must forget.

A mind cannot retain every experience. An institution cannot preserve every event. A scientific theory cannot represent every detail of every observation. Even the observable universe contains only a fragment of the information contained in its own history.

This fact is so familiar that it is easy to overlook its significance. Every act of remembering presupposes forgetting. Every representation presupposes compression. Every record of the past is already a selection from a history too large to preserve in full.

The central question of this book emerges from that observation.

If the past is compressed, when can it still be recovered?

More precisely: under what conditions can a system reconstruct answers to questions about alternative histories from a representation that does not contain those histories explicitly?

This question appears in many forms. Machine learning systems must answer questions about training data that no longer exists at inference time. Scientific theories must account for observations that vastly exceed what any single model can represent. Legal institutions must reason about social realities from archives that preserve only fragments of them. Cosmologists must infer the history of the universe from signals that survive in the present configuration of matter, radiation, and geometry.

Although these domains appear unrelated, they share a common structure. Each begins with a space of possible histories---a space of trajectories through which the system might have passed. Each applies a compression that preserves some distinctions and discards others. Each attempts to answer questions about trajectories that are no longer directly accessible.

The framework developed in this book studies that shared structure.

Part~One develops the mathematics of reconstruction from compressed representations. Part~Two examines three domains---learning systems, semantic inference, and physical field dynamics---in which reconstruction can be proved or rigorously tested under explicit conditions. Part~Three extracts a diagnostic methodology applicable to any domain where reconstruction is claimed or required. Part~Four applies that methodology to memory, institutions, and cosmology. Part~Five asks what these results imply about explanation, representation, and the geometry of alternative histories.

The goal throughout is not to prove that reconstruction is always possible. In many systems it is not, and the framework is designed to identify precisely when and why it fails. Nor is the goal to defend any particular theory of memory, cognition, cosmological dynamics, or institutional design. The goal is narrower and, I believe, more fundamental: to identify the conditions under which reconstruction from compressed representations becomes possible at all, and to characterize precisely what happens when those conditions are not met.

Those conditions turn out to be three. The history's influence on the query must decay with interaction order---the stability condition. The compression must preserve the information the query requires---the sufficiency condition. And the compression must be calibrated to the stability structure so that both conditions are controlled by a single depth parameter---the alignment condition. When all three hold, the present serves as a recoverable record of the past. When any one fails, it fails in a specific, diagnosable way that the framework can characterize precisely.

The chapters that follow develop this claim from its mathematical foundations to its philosophical consequences.

They begin with compression.

They end with memory.

Between those two ideas lies a single question: what kind of world allows its history to remain recoverable after it has been compressed?

This book is an attempt to answer that question.
